\chapter*{预备知识、记号、回顾与补充}
\addcontentsline{toc}{chapter}{预备知识、记号、回顾与补充}
\label{chap:ag-preliminaries}
\paragraph{预备知识}
本课程假定读者熟悉\newterm{ag-scheme}{概形}{scheme}的语言，程度相当于 Hartshorne 著作第二章截至第 II.5 节的内容. 因而，读者应充分理解层，特别是环层与模层，以及交换代数中张量积、$\operatorname{Hom}$ 等基本运算的层化. 概形的大多数``局部''性质依赖于有时相当困难、甚至非常困难的交换代数. 本文不再回顾局部化、平坦性、Krull 维数，以及 Noether 环、Artin 环、局部环、赋值环、离散赋值环和唯一分解整环等概念. 读者最好已经接触过整闭环（正规环）与正则环；我们会对这两类环作一些回顾.
\paragraph{记号}
回顾，概形是局部环化空间 $(X,\mathcal O_X)$，它局部同构于仿射概形 $(\operatorname{Spec}(A),\mathcal O_{\operatorname{Spec}(A)})$. 因而，拓扑空间 $X$ 具有``谱''的性质，即每个不可约闭集都有一个泛点. 环层 $\mathcal O_X$（直观上是``$X$ 上正则函数''所成的层）称为\newterm{ag-structure-sheaf}{结构层}{structure sheaf}. $\mathcal O_X$-模层简称为 $\mathcal O_X$-模. $\mathcal O_X$ 在点 $x\in X$ 处的茎是一个局部环，记为 $\mathcal O_{X,x}$，其极大理想记为 $\mathfrak m_{X,x}$，剩余域记为 $\kappa(x)$. 为简化记号，概形简记为 $X$，仿射概形简记为 $\operatorname{Spec}(A)$. 回顾，概形范畴中存在纤维积，但 $X\times_ZY$ 的底拓扑空间并不是底拓扑空间的纤维积.

概形态射 $Y\longrightarrow X$ 是局部环化空间的态射，严格说应始终记为 $(f,f^\sharp)$，其中 $f$ 是底拓扑空间之间的连续映射，而 $f^\sharp:f^{-1}\mathcal O_X\longrightarrow\mathcal O_Y$ 是 $Y$ 上环层的态射，满足：对每个 $y\in Y$，都有 $(f^\sharp)^{-1}(\mathfrak m_{Y,y})=\mathfrak m_{X,f(y)}$. 不过，我们仍简写为 $f$. 要区分将 $X$ 上的层送到 $Y$ 上的层的函子 $f^{-1}$，以及将 $\mathcal O_X$-模送到 $\mathcal O_Y$-模的函子 $f^*$. 当然，给定一个 $\mathcal O_X$-模 $M$，有
\begin{align*}
f^*(M)=f^{-1}(M)\otimes_{f^{-1}(\mathcal O_X)}\mathcal O_Y,
\end{align*}
其中张量积应在层的意义下理解.
\paragraph{回顾：凝聚层与拟凝聚层}
设 $(X,\mathcal O_X)$ 是局部环化空间. $X$ 上的\newterm{ag-quasi-coherent}{拟凝聚层}{quasi-coherent sheaf}，亦称拟凝聚 $\mathcal O_X$-模，是满足下列条件的 $\mathcal O_X$-模 $\mathcal F$：每个点 $x\in X$ 都有邻域 $U$，使得在 $U$ 上存在 $\mathcal O_U$-模的正合列
\begin{align*}
\mathcal O_U^{(I)}\longrightarrow\mathcal O_U^{(J)}\twoheadrightarrow\mathcal F|_U,
\end{align*}
其中 $I,J$ 是集合，$\mathcal O_U^{(I)}$ 表示层的直和 $\bigoplus_{i\in I}\mathcal O_U$. 若可取 $I$ 为空集（此时 $\mathcal O_U^{(J)}\xrightarrow{\sim}\mathcal F|_U$），则称 $\mathcal F$ \newterm{ag-locally-free}{局部自由}{locally free}. 若可取 $J$ 为有限集，则称 $\mathcal F$ 为\newterm{ag-finite-type}{有限型}{of finite type}. 在此情形下，若还可取 $I$ 为有限集（必要时缩小邻域 $U$），则称 $\mathcal F$ 为\newterm{ag-finite-presentation}{有限表现}{of finite presentation}. 最后，若对每个开集 $U$，$\mathcal F|_U$ 的每个有限型子层都有限表现，则称 $\mathcal F$ 为\newterm{ag-coherent}{凝聚层}{coherent sheaf}. 以 $\operatorname{QCoh}(X)$、$\operatorname{Coh}(X)$ 分别记 $\mathcal O_X\text{-}\operatorname{Mod}$ 中由 $X$ 上拟凝聚层、凝聚层组成的满子范畴. 范畴 $\operatorname{Coh}(X)$ 对核、余核和扩张封闭，特别地，它是 Abel 范畴. 对任意态射 $f:(Y,\mathcal O_Y)\longrightarrow(X,\mathcal O_X)$，函子 $f^*$ 将 $\operatorname{QCoh}(X)$ 映入 $\operatorname{QCoh}(Y)$，并将 $\operatorname{Coh}(X)$ 映入 $\operatorname{Coh}(Y)$. 注意，在这样一般的情形下，$\mathcal O_X$ 显然拟凝聚，却未必凝聚！若 $\mathcal O_X$ 凝聚，则 $\mathcal O_X$-模 $\mathcal F$ 凝聚当且仅当它有限表现.

\begin{example*}
当 $X$ 是单点时，就得到通常的（交换）代数概念. 环 $\overline{\mathbb Z}_p$（$\mathbb Z_p$ 在 $\mathbb Q_p$ 的一个代数闭包中的整闭包）是一个虽非 Noether 环却仍很有用的凝聚环例子. 环 $k[X_i]_{i\in\mathbb N}$ 是另一个例子.
\end{example*}
对于概形，这些范畴有良好的性质. 若 $X$ 是概形，则 $\operatorname{QCoh}(X)$ 对核、余核和扩张封闭，特别地，它是 Abel 范畴. 对分离且拟紧的态射 $f:Y\longrightarrow X$，直接像函子 $f_*$ 将 $\operatorname{QCoh}(Y)$ 映入 $\operatorname{QCoh}(X)$，参见 [EGA1, 9.2]. 若 $X$ 还是 Noether 概形，则 $\mathcal O_X$ 凝聚，并且 $\operatorname{Coh}(X)$ 是 $\mathcal O_X\text{-}\operatorname{Mod}$ 中包含所有有限秩局部自由层、且对直和、核、余核和扩张封闭的最小本质满子范畴.

这些性质来自仿射情形下的下列基本结果：若 $X=\operatorname{Spec}(A)$，则函子 $\mathcal F\mapsto\Gamma(X,\mathcal F)$ 诱导 $\operatorname{QCoh}(X)$ 与 $A\text{-}\operatorname{Mod}$ 的范畴等价. 若 $A$ 还是 Noether 环，则这一等价将凝聚层送到有限生成模. 一个拟逆函子可如下给出：对每个 $A$-模 $M$，令 $\widetilde M$ 为预层 $U\mapsto M\otimes_A\mathcal O_{\operatorname{Spec}(A)}(U)$ 的伴随层. 若开集 $U$ 是``主开集''，即对某个 $f\in A$ 有 $U=D_f$，则可以证明
\begin{align*}
\Gamma(U,\widetilde M)&=M\otimes_A\mathcal O_{\operatorname{Spec}(A)}(U)\\
&=M\otimes_A A_{(f)}=M_{(f)},
\end{align*}
其中 $M_{(f)}$ 是 $M$ 关于由 $f$ 生成的乘性子集的局部化. 参见 [Hartshorne II.5].

由于（拟）凝聚性是局部性质，概形 $X$ 上的层 $\mathcal F$（拟）凝聚当且仅当存在仿射覆盖 $X=\bigcup_iU_i$，使得 $\mathcal F|_{U_i}$ 在 $U_i$ 上拟凝聚. 对于具有显式仿射覆盖的概形，有时可以由仿射情形得到凝聚层的描述. 例如，在 $\mathbb P_k^1=\operatorname{Spec}(k[T])\cup\operatorname{Spec}(k[T^{-1}])$ 上，一个凝聚层由一个 $K[T]$-模 $M_+$、一个 $K[T^{-1}]$-模 $M_-$，以及一组粘合资料给出，即 $K[T,T^{-1}]$-模的同构
\begin{align*}
M_+\otimes_{K[T]}K[T,T^{-1}]\xrightarrow{\sim}M_-\otimes_{K[T^{-1}]}K[T,T^{-1}].
\end{align*}
对于射影簇，更一般地，对于分次环的齐次谱，我们稍后将看到，可以用分次模来描述凝聚模.
